\documentclass[12pt,a4paper]{article}
\usepackage[utf8]{inputenc}
\usepackage[T1]{fontenc}
\usepackage{amsmath,amssymb,amsfonts}
\usepackage{mathtools}
\usepackage{booktabs}
\usepackage{array}
\usepackage{geometry}
\usepackage{hyperref}
\usepackage{cite}
\geometry{margin=2.5cm}

\newcommand{\MeV}{\,\text{MeV}}
\newcommand{\GeV}{\,\text{GeV}}
\newcommand{\fpi}{f_\pi}
\newcommand{\qqbar}{\langle\bar{q}q\rangle}
\newcommand{\zch}{\mathfrak{z}}

\title{The Koide formula and quark-lepton mass relations:\\
evidence and statistical assessment}

\author{A.~Rivero\thanks{Universidad de Zaragoza.
Manuscript prepared with AI assistance (Claude/Anthropic).}}
\date{March 2026}

\begin{document}
\maketitle

\begin{abstract}
We compile and statistically evaluate a family of mass relations among Standard Model
fermions built on the Koide ratio $Q(m_1,m_2,m_3)\equiv(m_1+m_2+m_3)/(\sqrt{m_1}+\sqrt{m_2}+\sqrt{m_3})^2$.
The original charged-lepton relation $Q(m_e,m_\mu,m_\tau)=2/3$ holds to $0.001\%$
($0.91\sigma$ in $m_\tau$); a Monte Carlo look-elsewhere analysis over all fermion triples
gives $2.8\sigma$.  Extending the formula to signed mass charges
$\zch_k=\pm\sqrt{m_k}$, we find the quark triples
$(-m_s,m_c,m_b)$ and $(m_c,m_b,m_t)$ satisfy $Q\approx 2/3$ to $1.2\%$ and $0.4\%$
respectively, both obligatorily mixing up- and down-type quarks.
A seed triple $(0,m_s,m_c)$ with the O'Raifeartaigh mass ratio
$\sqrt{m_s}/\sqrt{m_c}=2-\sqrt{3}$ is confirmed to $1.2\%$.
The vacuum-charge doubling relation $\sqrt{m_b}=3\sqrt{m_s}+\sqrt{m_c}$ predicts
$m_b=4177$ MeV ($0.1\sigma$ from PDG).
The Koide phase parameter satisfies
$\delta_0\bmod 2\pi/3 \approx 2/9$ to 33 ppm ($3.1\sigma$ with look-elsewhere, $4.5\sigma$ specific).
A meson triple $(-\pi,D_s,B)$ extends $Q\approx 2/3$ to composite hadrons at $0.1\%$.
The dual Koide $Q(1/m_d,1/m_s,1/m_b)=0.665$ ($0.22\%$ from $2/3$) suggests a
Seiberg-dual structure.
The de Vries angle $R=(\sqrt{19}-3)(\sqrt{19}-\sqrt{3})/16=0.22310\ldots$ matches
the on-shell $\sin^2\theta_W$ to $0.13\sigma$,
and the Gatto--Sartori--Tonin relation $\sin\theta_C=\sqrt{m_d/m_s}$ holds to $0.3\%$.
Taken individually, several of these relations have modest significance; taken together,
their interconnections---shared algebraic structure, overlapping mass triplets, and
common Koide phase---point to a single organizing principle.
The companion paper~[I] constructs the $\mathcal{N}=1$ Lagrangian that realizes this
principle via an O'Raifeartaigh superpotential with structural invariant $Q=2/3$.
\end{abstract}

%======================================================================
\section{Introduction}
%======================================================================

\subsection{The Koide formula}

In 1981, Koide derived a relation among the charged lepton masses from a composite model
with $S_3$ permutation symmetry acting on preon constituents~\cite{Koide1983}.  Stripped of its
model-dependent origin, the relation states that the ratio
\begin{equation}
\boxed{Q(m_1,m_2,m_3) \;\equiv\; \frac{m_1+m_2+m_3}{\bigl(\sqrt{m_1}+\sqrt{m_2}+\sqrt{m_3}\bigr)^2}}
\end{equation}
equals $2/3$ when evaluated on the electron, muon, and tau masses.  With current PDG
values~\cite{PDG2024},
\begin{equation}
Q(m_e,\,m_\mu,\,m_\tau) \;=\; 0.666661 \pm 0.000007\,,
\end{equation}
where the uncertainty is dominated by $\delta m_\tau = 0.12$ MeV.  If the formula is
exact, it predicts $m_\tau = 1776.97$ MeV, consistent with the experimental value
$1776.86 \pm 0.12$ MeV to $0.91\sigma$.

Foot reformulated this as a geometric condition: the vector of square roots of masses
makes a $45^\circ$ angle with the democratic vector $(1,1,1)$~\cite{Foot1994}.  Equivalently, in the
angular parametrization
\begin{equation}
\sqrt{m_k} \;=\; \sqrt{M_0}\,\Bigl(1 + \sqrt{2}\,\cos\!\Bigl(\tfrac{2\pi k}{3}+\delta_0\Bigr)\Bigr)\,,
\qquad k=0,1,2\,,
\end{equation}
the condition $Q=2/3$ is satisfied identically for any $M_0$ and $\delta_0$, and the
three masses are determined by two parameters.  The phase for charged leptons is
$\delta_0 = 0.2222$ (in natural units of $2\pi/3$).

The precision of the lepton relation---better than one part in $10^5$---has prompted
four decades of theoretical attempts to explain it, ranging from discrete symmetries
($S_3$, $A_4$) acting on flavon fields~\cite{Koide1990,Ma2006} to radiative mechanisms that protect
the formula against renormalization group running~\cite{XZZ2008}.  Yet no consensus model exists.
The formula has been called an ``unexplained coincidence'' and a ``tantalizing hint''
in roughly equal measure.

\subsection{The skeptic's objection: look-elsewhere and scheme dependence}

Two objections are standard.

\paragraph{Look-elsewhere.}  With six quarks and three charged leptons forming
$\binom{9}{3} = 84$ triples (or more, allowing signed square roots), a near-hit at
$Q=2/3$ somewhere in the data is not implausible by chance.  We address this head-on.
A Monte Carlo simulation drawing nine masses log-uniformly over the range
$[0.1,\,200{,}000]$ MeV and scanning all 84 triples times 8 sign assignments finds that
a match as precise as the charged leptons occurs in only $0.24\%$ of random mass
spectra.  This is a look-elsewhere corrected significance of $2.8\sigma$.  Not
overwhelming---but the lepton Koide is only the first entry in a longer list.

\paragraph{Renormalization scheme.}  The Koide formula is stated in terms of pole masses (or
$\overline{\text{MS}}$ masses evaluated at the mass itself, $m_q(m_q)$).  Running to the
GUT scale shifts $Q$ by a few percent for quarks and by less than $0.1\%$ for leptons~\cite{XZZ2008,Esposito1995}.  This is sometimes used to argue that the relation is an accident of the
infrared.  But it is equally true that many well-established mass relations---the
Gell-Mann--Okubo formula, the Gatto--Sartori--Tonin (GST) relation---are statements about
low-energy masses, not running couplings at a unification scale.  If the Koide relation
originates from boundary conditions on an effective superpotential (as argued in~[I]),
then the infrared is exactly where it should hold.

\subsection{Beyond the leptons: a web of mass relations}

The real case for the Koide formula rests not on a single relation but on a web of
interlocking observations.  In 2011, Rodejohann and Zhang~\cite{RodejohannZhang2011} noted that the triple
$(m_c, m_b, m_t)$ satisfies $Q = 0.6693$, within $0.4\%$ of $2/3$.  This was surprising
because it mixes quark generations and charge sectors.  In~\cite{Rivero2012}, one of us showed that
allowing signed square roots reveals a further triple $(-m_s, m_c, m_b)$ with $Q = 0.675$
($1.2\%$ from $2/3$), and that these triples chain together: the output of one is the
input of the next, covering five of the six quark masses.

The present paper collects all known mass relations of this type---ten observations
spanning charged leptons, quarks, and pseudoscalar mesons---and evaluates each one
statistically.  The observations, in order of discovery, are:

\begin{enumerate}
\item \textbf{The charged-lepton Koide}~\cite{Koide1983}:
   $Q(m_e,m_\mu,m_\tau) = 2/3$ to $0.001\%$.  Significance: $0.91\sigma$ in $m_\tau$;
   $2.8\sigma$ after look-elsewhere.

\item \textbf{The quark seed}~\cite{Rivero2012}:
   $Q(0,\,m_s,\,m_c) \approx 2/3$, meaning $\sqrt{m_s}/\sqrt{m_c} \approx 2-\sqrt{3}$.
   Observed ratio $0.271$, theoretical $0.268$; deviation $1.2\%$.

\item \textbf{The signed quark triple}~\cite{Rivero2012}:
   $Q(-m_s,\,m_c,\,m_b) = 0.675$, deviation $1.24\%$ from $2/3$.

\item \textbf{The heavy triple}~\cite{RodejohannZhang2011}:
   $Q(m_c,\,m_b,\,m_t) = 0.6693$, deviation $0.40\%$.

\item \textbf{The meson triple}~\cite{Rivero2007}:
   $Q(-m_\pi,\,m_{D_s},\,m_B) = 0.6674$, deviation $0.10\%$.

\item \textbf{$v_0$-doubling} (this work):
   $\sqrt{m_b} = 3\sqrt{m_s} + \sqrt{m_c}$ to $0.07\%$, predicting $m_b = 4177$ MeV
   (PDG: $4180 \pm 30$ MeV, $0.1\sigma$).

\item \textbf{The phase relation}~\cite{Zenczykowski2012}:
   $\delta_0 \bmod 2\pi/3 \approx 2/9$ to 33 ppm.  Significance: $3.1\sigma$ with
   look-elsewhere over 128 simple fractions; $4.5\sigma$ for the specific value $2/9$.

\item \textbf{The dual Koide} (this work):
   $Q(1/m_d,\,1/m_s,\,1/m_b) = 0.665$, $0.22\%$ from $2/3$.
   Suggests a Seiberg-dual seesaw structure $M_j \propto 1/m_j$.

\item \textbf{The cyclic echo} (this work):
   Iterating the Koide equation as $t \to c \to s \to \text{stall} \to c$ produces an
   algebraically exact cycle (agreement to machine epsilon).

\item \textbf{The second scaling fails} (this work):
    $M_2 = 9M_0$, $\delta_2 = 9\delta_0$ does \emph{not} produce $(c,b,t)$ or any physical
    triple.  The pattern stops at one step.  This negative result constrains the
    mechanism.
\end{enumerate}

Two additional relations, while not Koide triples, belong to the same family:

\begin{itemize}
\item \textbf{The de Vries angle}~\cite{Rivero2005b}:
  $R = (\sqrt{19}-3)(\sqrt{19}-\sqrt{3})/16 = 0.22310\ldots$ matches
  $\sin^2\theta_W(\text{on-shell})$ to $0.13\sigma$ in $M_W$.

\item \textbf{The GST/Oakes relation}~\cite{GST1968}:
  $\sin\theta_C = \sqrt{m_d/m_s} = 0.224$, matching $|V_{us}| = 0.2243$ to $0.3\%$.
\end{itemize}

\subsection{What this paper adds}

Previous discussions of these relations appear scattered across individual papers,
conference proceedings, and blog posts.  No systematic statistical assessment exists.
This paper provides three things:

\textbf{First}, a self-contained compilation of all the observations, with current PDG values
and error bars, in a single reference.  We state every number with its uncertainty and
pull, not just a percentage deviation.

\textbf{Second}, a rigorous look-elsewhere analysis.  For each observation we define the trial
space---how many similar relations could have been examined---and compute corrected
$p$-values.  We are explicit about which observations are significant in isolation
($2.8\sigma$ for the lepton Koide, $3.1\sigma$ for the phase relation) and which derive
their significance only from the structural context ($v_0$-doubling, dual Koide).

\textbf{Third}, an argument that the observations are not independent.  The quark seed, the
signed triple, the heavy triple, and the $v_0$-doubling relation share three of the same
masses ($m_s$, $m_c$, $m_b$) and are linked by a common algebraic structure.
An exhaustive scan of all $\binom{6}{3} = 20$ quark triples reveals that \emph{every} triple
close to $Q = 2/3$ obligatorily mixes up-type and down-type quarks~\cite{Rivero2012}; pure
same-charge triples $(d,s,b)$ and $(u,c,t)$ give $Q = 0.73$ and $Q = 0.85$
respectively, both far from $2/3$.  This is not a selection effect: while $90\%$ of
random triples from six quarks are mixed by combinatorics, the $100\%$ rate among
Koide-satisfying triples is a structural feature pointing toward CKM-type mixing.

The companion paper~[I] constructs the $\mathcal{N}=1$ supersymmetric Lagrangian that
produces these relations.  The present paper is deliberately agnostic about the Lagrangian.
We state the empirical relations, assess their statistical weight, and leave the dynamical
interpretation to~[I].
The question this paper asks is precise: \emph{given the data, is there evidence for a common
algebraic structure in the fermion mass spectrum, or can the observations be dismissed as
look-elsewhere artifacts?}

\subsection{Conventions and mass inputs}

Throughout this paper, quark masses refer to $\overline{\text{MS}}$ values at the
self-energy scale $m_q(m_q)$ for $c$ and $b$, and the pole mass for the top quark:
$m_s = 93.4 \pm 5.8$ MeV, $m_c = 1270 \pm 20$ MeV, $m_b = 4180 \pm 30$ MeV,
$m_t = 172.76 \pm 0.30$ GeV~\cite{PDG2024}.  Light quark masses are $m_u = 2.16 \pm 0.49$ MeV,
$m_d = 4.67 \pm 0.48$ MeV.  Lepton masses are pole masses: $m_e = 0.51100$ MeV,
$m_\mu = 105.658$ MeV, $m_\tau = 1776.86 \pm 0.12$ MeV.  Meson masses are PDG
averages: $m_\pi = 139.57$ MeV, $m_{D_s} = 1968.35$ MeV, $m_B = 5279.34$ MeV.

We use natural units $c = \hbar = 1$.  The Koide ratio is denoted $Q$.  We write
$\zch_k = \pm\sqrt{m_k}$ for signed mass charges, with the sign conventions
specified for each triple.  Pulls are quoted as
$(x_{\text{pred}} - x_{\text{obs}})/\sigma_{\text{obs}}$ in units of $\sigma$.

\subsection{Outline}

Section~2 reviews the Koide parametrization and the seed structure.
Section~3 presents the quark chain: the seed, the signed triple, and the heavy triple.
Section~4 discusses the $v_0$-doubling and the overlap prediction.
Section~5 extends the analysis to pseudoscalar mesons.
Section~6 covers the heavy triple $(c,b,t)$.
Section~7 presents mass relations and iterative chains.
Section~8 discusses the de Vries angle and the Casimir equation.
Section~9 gives the full statistical assessment.
Section~10 summarizes the conclusions and the connection to~[I].

%======================================================================
\section{The Koide parametrization and seed structure}
%======================================================================

\subsection{Charges, not masses}

The fundamental variable is the signed mass charge $\zch = z_0 + z_i
\in \mathbb{R}$, with the physical mass $m = \zch^2 \geq 0$.

\subsection{Derivation of $Q = 2/3$}

For three charges $\zch_k = z_0 + z_k$ with $z_0$ the mean
($\sum z_k = 0$):
\begin{equation}
Q = \frac{\sum \zch_k^2}{(\sum \zch_k)^2}
= \frac{3z_0^2 + \sum z_k^2}{9z_0^2}
= \frac{1}{3} + \frac{\langle z_k^2\rangle}{3z_0^2}\,.
\end{equation}
Therefore:
\begin{equation}
\boxed{Q = \frac{2}{3} \quad\Longleftrightarrow\quad
\langle z_k^2\rangle = z_0^2\,.}
\end{equation}

The two conditions are:
\begin{enumerate}
\item \textbf{Zero-sum perturbations:} $z_1 + z_2 + z_3 = 0$
(definition of $z_0$ as the mean charge).
\item \textbf{Energy balance:} The mean energy of the perturbations
equals the energy of the unperturbed charge:
$\langle z_k^2\rangle = z_0^2$.
\end{enumerate}

Condition~2 is the physics. It says the splitting scale matches the
base scale---the coefficient of variation equals unity:
$\sigma(\zch)/\bar{\zch} = 1$.
(Dimensionally, $\zch$ has units $\MeV^{1/2}$, so
$z_0^2$ has units MeV---a mass scale, not an energy density.
The ``energy balance'' terminology is metaphorical.) This is neither perturbative
($\sigma \ll \bar{\zch}$, giving $Q \to 1/3$, all equal) nor
dominant ($\sigma \gg \bar{\zch}$, giving $Q \to 1$, one mass
dominates). The Koide value $2/3$ is the equipartition point.

\subsection{Why one mass can vanish}

If $\langle z_k^2\rangle = z_0^2$, a perturbation $z_k$ can reach
$|z_k| \sim z_0$, making $\zch_k = z_0 + z_k \approx 0$. A vanishing
mass is the \emph{natural extreme} of energy balance, not fine-tuning.

\subsection{Seed triples: the $(0, a, b)$ structure}

For a triple $(0, \zch_a, \zch_b)$ with one massless member, the
Koide ratio becomes:
\begin{equation}
Q = \frac{\zch_a^2 + \zch_b^2}{(\zch_a + \zch_b)^2}\,.
\end{equation}
Setting $Q = 2/3$ and writing $r = \zch_a/\zch_b$ (with $r < 1$):
\begin{equation}
3(r^2 + 1) = 2(r + 1)^2 \quad\Longrightarrow\quad r^2 - 4r + 1 = 0\,,
\end{equation}
giving:
\begin{equation}
\boxed{r = 2 - \sqrt{3} \approx 0.26795\,.}
\end{equation}

Equivalently, in mass space:
$m_{\text{light}}/m_{\text{heavy}} = (2 - \sqrt{3})^2 = 7 - 4\sqrt{3}
\approx 0.07180$.

This is a \emph{universal} prediction: any Koide triple with one
massless member has the same charge ratio between the two massive
members, regardless of their absolute scale.

\subsection{The quark seed: $(0, s, c)$}

The strange and charm quark charges give:
\begin{equation}
\frac{\sqrt{m_s}}{\sqrt{m_c}} = \frac{9.67}{35.64} = 0.2712\,,
\qquad r_{\text{Koide}} = 0.2679\,, \qquad \text{deviation: } 1.2\%\,.
\end{equation}

This triple has $Q = 0.6644$, deviating from $2/3$ by only $0.35\%$.
Its Koide scale is $v_0^2 = (\zch_s + \zch_c)^2/9 \approx 228\MeV$.

\paragraph{Scale ambiguity in the quark seed.}
Throughout this paper, quark masses follow PDG conventions:
$m_s(2\GeV) = 93.4\MeV$, $m_c(m_c) = 1.27\GeV$,
$m_b(m_b) = 4.18\GeV$, $m_t = 172.76\GeV$ (pole mass). This is a
mixed convention, with the light quark evaluated at 2~GeV and heavier
quarks at their own scale. QCD mass ratios are RG invariants at
leading order (the anomalous dimension $\gamma_m$ is flavor-universal),
so $Q$ is scheme-independent. However, the seed ratio
$\sqrt{m_s}/\sqrt{m_c}$ uses masses at \emph{different} scales:
at a common $\overline{\text{MS}}$ scale the ratio shifts to
$\sim 0.29$, worsening the seed match from 1.2\% to $\sim 8\%$.
This is a significant limitation of the quark seed claim. The meson
seed, which uses unambiguous physical (pole) masses, is free of this
scheme issue and provides the cleaner test of the seed hypothesis.

A historical note: Harari, Haut, and Weyers~\cite{Harari} observed
a quark mass ratio pattern that, in the language of the present paper,
corresponds to an approximate seed triple---$(0, d, s)$---at a time
when the up quark mass was
still thought to be possibly zero. With current masses,
$\sqrt{m_d}/\sqrt{m_s} = 0.224$, deviating from $2 - \sqrt{3}$ by
$16\%$---too large for a good Koide triple. But if $m_s$ were
$\sim 150\MeV$ and $m_d \sim 7\MeV$ (values common in that era), the
ratio improves considerably. The seed idea was already present, ahead
of its time.

\subsection{The meson seed: $(0, \pi, D_s)$}

The pion and $D_s$ meson charges give:
\begin{equation}
\frac{\sqrt{m_\pi}}{\sqrt{m_{D_s}}} = \frac{11.81}{44.37} = 0.2663\,,
\qquad r_{\text{Koide}} = 0.2679\,, \qquad \text{deviation: } 0.6\%\,.
\end{equation}

This gives $Q = 0.6679$, deviating by $0.18\%$. It is the
\textbf{meson counterpart} of the quark seed $(0, s, c)$---both
involve the strange-charm sector. Indeed, $D_s = c\bar{s}$ is the
meson that directly contains the $s$ and $c$ quarks.

\subsection{The lepton triple as an ``almost-seed''}

The charged lepton triple $(e, \mu, \tau)$ has $\zch_e = 0.715\MeV^{1/2}$,
which is small but nonzero. If $m_e$ were zero, we would need
$\sqrt{m_\mu}/\sqrt{m_\tau} = 2 - \sqrt{3} = 0.2679$; the actual value
is $0.2439$, deviating by $9\%$. So the leptons are \emph{not} an
exact seed---the electron, while very light, is not massless. But the
triple is clearly the same family: the near-critical Koide phase puts
one member close to zero.

\subsection{Scale coincidence}

The $v_0^2$ values of the seed triples cluster remarkably:

\begin{center}
\begin{tabular}{lcc}
\toprule
Triple & $v_0^2$ (MeV) & $v_0$ ($\MeV^{1/2}$) \\
\midrule
$(e, \mu, \tau)$ & 314 & 17.72 \\
$(0, \pi, D_s)$ & 351 & 18.73 \\
$(0, s, c)$ & 228 & 15.10 \\
\bottomrule
\end{tabular}
\end{center}

The lepton and meson seed triples share a Koide scale around
310--350~MeV, squarely in the QCD scale range. This is the
sBootstrap connection made quantitative: \textbf{the lepton triple
and the meson seed triple operate at the same scale.}

For the lepton triple, the coincidence is numerically sharp:
\begin{equation}
M_0 \;\approx\; \frac{\fpi^2}{27}\,,
\end{equation}
to $0.3\%$: the right-hand side gives $314.9\MeV$ versus
the fitted $M_0 = 313.8\MeV$.  Equivalently,
$v_0 \approx \fpi/(3\sqrt{3})$.  If this identification is exact
and $\delta_0 \bmod 2\pi/3 = 2/9$ (the three-instanton value),
all three charged lepton masses follow from~$\fpi$ alone,
each predicted to~${\sim}0.3\%$.

%======================================================================
\section{The quark chain}
%======================================================================

\subsection{From seed to full triple}

The seed triples grow into full triples upon symmetry breaking:

\begin{center}
\begin{tabular}{lccc}
\toprule
Sector & Seed triple & Full triple & $v_0^2$ (full) \\
\midrule
Quarks & $(0, s, c)$ & $(-s, c, b)$ & 913 MeV \\
Mesons & $(0, \pi, D_s)$ & $(-\pi, D_s, B)$ & 1230 MeV \\
Leptons & $(0 \approx e, \mu, \tau)$ & $(e, \mu, \tau)$ & 314 MeV \\
\bottomrule
\end{tabular}
\end{center}

In each case:
\begin{itemize}
\item The seed has a massless (or near-massless) member, two massive
members in the ratio $2 - \sqrt{3}$, and $Q \approx 2/3$.
\item The full triple has \textbf{no} massless member: the zero has
been lifted (to $m_e$, or the sign-flipped strange/pion mass). The
$b$ quark and $B$ meson are always present as flavor degrees of
freedom; what the breaking does is give them their \emph{distinctive
masses}---the specific values that complete the Koide triple. Before
breaking, the $b$-direction sits at the flat direction ($\zch = 0$);
after, it acquires $\zch_b = 64.65\MeV^{1/2}$, the value selected
by the energy balance with $c$ and the sign-flipped $s$. The bloom
is not particle creation but \emph{mass generation}.
\end{itemize}

The ``tale of two tuples'' is this: \textbf{each sector presents
first as a seed with a massless member, then ``blooms'' into a
full triple when SUSY breaks.} The massless member is the necessary
initial condition; the breaking lifts it and generates the full spectrum.

\subsection{Precision summary}

\begin{center}
\begin{tabular}{lcccc}
\toprule
Triple & $Q$ & Dev.\ from $2/3$ & $v_0^2$ (MeV) & $\delta$ (rad) \\
\midrule
\multicolumn{5}{c}{\textit{Seed triples}} \\
$(0, \pi, D_s)$ & 0.6679 & 0.18\% & 351 & $\approx 3\pi/4$ \\
$(0, s, c)$ & 0.6644 & 0.35\% & 228 & $\approx 3\pi/4$ \\
$(0 \approx e, \mu, \tau)$ & 0.6667 & 0.001\% & 314 & 2.317 \\
\midrule
\multicolumn{5}{c}{\textit{Full triples}} \\
$(e, \mu, \tau)$ & 0.66666 & 0.001\% & 314 & 2.317 \\
$(-\pi, D_s, B)$ & 0.6674 & 0.10\% & 1230 & 2.806 \\
$(-s, c, b)$ & 0.6750 & 1.24\% & 913 & 2.744 \\
$(c, b, t)$ & 0.6695 & 0.42\% & 29\,580 & 2.163 \\
\bottomrule
\end{tabular}
\end{center}

The Koide phases $\delta$ cluster in the range $2.2$--$2.8$~rad
($126^\circ$--$161^\circ$), all displaced from the seed value
$3\pi/4 \approx 2.356$. The lepton triple barely moved from its
seed; the quark and meson triples bloomed further. All phases
lie within one $\mathbb{Z}_3$ fundamental domain ($[0, 2\pi/3]$
modulo $2\pi/3$).

The lepton triple is unique: seed and full triple are nearly
identical, because the electron mass is small enough to serve as
the ``zero'' while being nonzero enough to make $Q$ extremely
precise. The $(-s,c,b)$ triple has the worst deviation (1.24\%);
this is expected because the bloom is a nonperturbative process
that does not exactly preserve $Q$. The $v_0$-doubling constrains
$m_b$ independently (to 0.07\% via $\sqrt{m_b} = 3\sqrt{m_s} +
\sqrt{m_c}$), so the 1.24\% Koide deviation is the residual after
the more fundamental $v_0$-constraint is satisfied.

\subsection{Orthogonality of seed and full triples}

Each Koide triple defines a perturbation vector
$\vec{z} = (z_1, z_2, z_3)$ constrained to the plane $\sum z_k = 0$
with length $|\vec{z}|^2 = 3z_0^2$. The angle between seed and
full triples measures how much the breaking rotates the spectrum
in charge space:

\begin{center}
\begin{tabular}{lc}
\toprule
Pair of triples & Angle \\
\midrule
$(0,s,c)$ vs.\ $(-s,c,b)$ & $22^\circ$ \\
$(0,\pi,D_s)$ vs.\ $(-\pi,D_s,B)$ & $26^\circ$ \\
$(e,\mu,\tau)$ vs.\ $(0,\mu,\tau)$ & $0.8^\circ$ \\
$(0,s,c)$ vs.\ $(0,\pi,D_s)$ & $0.3^\circ$ \\
\bottomrule
\end{tabular}
\end{center}

Two results are notable. First, the quark seed $(0,s,c)$ and meson
seed $(0,\pi,D_s)$ are \textbf{nearly parallel} ($0.3^\circ$): the
meson seed is the composite image of the quark seed, with the GOR
map preserving direction while changing scale. Second, the lepton
triple is essentially its own seed ($0.8^\circ$)---the bloom has
barely begun.

%======================================================================
\section{$v_0$-doubling and the overlap prediction}
%======================================================================

\subsection{The vacuum-charge doubling}

The vacuum charge $v_0 = \tfrac{1}{3}\sum\zch_k$ of the quark seed
is $v_0^{(s)} = (\sqrt{m_s} + \sqrt{m_c})/3 = 15.10\MeV^{1/2}$.
For the full triple, $v_0^{(f)} = (-\sqrt{m_s} + \sqrt{m_c} +
\sqrt{m_b})/3 = 30.21\MeV^{1/2}$. Their ratio is
\begin{equation}
  v_0^{(f)}/v_0^{(s)} = 2.0005 \approx 2\,.
\end{equation}
If the bloom exactly doubles the vacuum charge, then
$\sqrt{m_b} = 3\sqrt{m_s} + \sqrt{m_c}$, giving
$m_b = 4{,}177\MeV$---matching the PDG value
$4{,}180 \pm 30\MeV$ to $0.1\sigma$ ($0.06\sigma$ including
propagated input uncertainties). This relation is
independent of the Koide condition: it constrains $m_b$ from
$(m_s, m_c)$ alone, and is more precise (0.07\%) than the
overlapping-triples prediction (0.5\%).

\subsection{Tension with the exact seed}

A tension exists:
the exact seed Koide gives $m_c = 1{,}301\MeV$ (not the PDG
$1{,}270$), and applying the $v_0$-doubling to this exact seed
value gives $m_b = 4{,}233\MeV$ ($1.8\sigma$ from PDG). The
$v_0$-doubling works \emph{with the physical} $m_c$, not the exact
seed value---so these two constraints are mutually inconsistent at the
stated precision. The resolution may be that the seed Koide fixes the
ratio $m_c/m_s$ only approximately (the 2.4\% charm deviation being a
higher-order correction to the $gv/m = \sqrt{3}$ condition), while the
$v_0$-doubling constrains $m_b$ directly from the physical masses. In
this reading, the seed condition is the approximate one and the
$v_0$-doubling the more fundamental. This hierarchy is not derived;
it is imposed by the data.

\subsection{The overlapping-triples prediction}

Requiring $Q = 2/3$ on both $(-s,c,b)$ and $(c,b,t)$ simultaneously
with inputs $(m_s, m_t)$ gives $m_b = 4159\MeV$ (0.5\% from PDG) and
$m_c = 1369\MeV$ (7.8\% from PDG). The $v_0$-doubling is more
precise for $m_b$ ($0.07\%$ vs.\ $0.5\%$), while the overlapping-triples
method additionally constrains $m_c$.

\subsection{The exact seed and what breaks it}

The seed configuration $(0, \zch_a, \zch_b)$ requires \emph{two}
independent conditions:
\begin{enumerate}
\item One member at $\zch = 0$ exactly (the flat direction).
\item The ratio of the massive members equals $2 - \sqrt{3}$:
\begin{equation}
\frac{\sqrt{m_a}}{\sqrt{m_b}} = 2 - \sqrt{3} \approx 0.26795\,.
\end{equation}
\end{enumerate}
These are logically distinct. The chiral limit ($m_u = m_d = 0$,
$\qqbar \neq 0$) enforces condition (1) for the meson sector: the
pion is an exact Goldstone boson, $m_\pi = 0$, and the
$B$-meson direction sits at
the flat direction before acquiring mass. But the chiral limit says
\emph{nothing} about condition (2)---it does not constrain
$\sqrt{m_{D_s}}/\sqrt{m_B}$ or $\sqrt{m_s}/\sqrt{m_c}$
to equal $2 - \sqrt{3}$.

The physical deviations from the exact ratio are small but
measurable:

\begin{center}
\begin{tabular}{lccc}
\toprule
Seed & Exact ratio & Physical ratio & Deviation \\
\midrule
$(0, \pi, D_s)$ & 0.2680 &
$\sqrt{m_\pi}/\sqrt{m_{D_s}} = 0.2663$ & $-0.6\%$ \\
$(0, s, c)$ & 0.2680 &
$\sqrt{m_s}/\sqrt{m_c} = 0.2712$ & $+1.2\%$ \\
$(0 \approx e, \mu, \tau)$ & 0.2680 &
$\sqrt{m_\mu}/\sqrt{m_\tau} = 0.2439$ & $-9.0\%$ \\
\bottomrule
\end{tabular}
\end{center}

The meson seed is closest to the exact ratio (0.6\% deviation).
The quark seed deviates by 1.2\%. The lepton ``seed'' has the worst ratio match ($-9\%$), which
is consistent with the observation that the lepton triple is
\emph{not} a true seed: $m_e \neq 0$, and the electron was never
at the flat direction. The leptons are an ``almost-seed'' whose
excellent $Q = 0.66666$ arises not from the seed structure but
from the smallness of $m_e$, which makes the two conditions
approximately degenerate.

%======================================================================
\section{The meson Koide triple: $(-\pi, D_s, B)$}
%======================================================================

With the pion charge negative, the triple $(-\pi, D_s, B)$ satisfies:
\begin{equation}
Q(-\pi, D_s, B) =
\frac{m_\pi + m_{D_s} + m_B}{(-\sqrt{m_\pi} + \sqrt{m_{D_s}} + \sqrt{m_B})^2}
= 0.6674\,,
\end{equation}
deviating from $2/3$ by $0.10\%$---the second most precise Koide
triple known. The energy balance:
\begin{equation}
\frac{\langle z_k^2\rangle}{z_0^2} = 1.002\,.
\end{equation}

For comparison, the ``standard'' all-positive quark triples $(d,s,b)$
and $(u,c,t)$ deviate from $Q = 2/3$ by 9.7\% and 27\% respectively
at $\overline{\text{MS}}$ masses. \textbf{The meson triple is more
precise than any quark triple}, including $(c,b,t)$ (0.42\%).

The pion's negative charge has physical meaning: as a pseudo-Goldstone
boson, the pion would be massless ($\zch = 0$) if chiral symmetry were
exact. It sits just below zero ($\zch_\pi = -11.81$), approaching zero
from the negative side as $m_q \to 0$. The sign of $\zch$ encodes
Goldstone nature: pseudo-Goldstone bosons sit on the negative side
of the critical point, having crossed through zero from positive
territory during chiral symmetry breaking.

\paragraph{Look-elsewhere correction.}
Eleven pseudoscalar mesons ($\pi$, $K$, $\eta$, $\eta'$, $D$,
$D_s$, $B$, $B_s$, $B_c$, $\eta_c$, $\eta_b$) yield
$\binom{11}{3} = 165$ triples, times
four independent sign patterns, giving 660 combinations.
Scanning all 660, only $(-\pi, D_s, B)$ and the closely related
$(-\pi, D_s, B_s)$ (deviation $0.14\%$) lie within $0.2\%$ of
$Q = 2/3$. A Monte Carlo with $10^7$ random triples drawn
from $[100,10{,}000]\MeV$ gives $p_{\text{single}} = 2.1 \times 10^{-4}$
for a match at $0.10\%$; correcting for 660 trials yields
$p_{\text{LEE}} = 0.13$ ($1.1\sigma$). The meson Koide triple
is therefore \emph{not} statistically significant on its own,
unlike the lepton triple ($2.8\sigma$ after LEE). Its interest
lies in the structural pattern: all six closest hits share the
sign assignment $(-\pi, D_{(s)}, B_{(s,c)})$, i.e.\ a negative
pion paired with charm-containing and bottom-containing mesons---the
combination predicted by the $SU(5)$ flavor structure.

\paragraph{Extended scan: gauge and Higgs bosons.}
Enlarging the scan to include the $W$ ($80.38\GeV$), $Z^0$ ($91.19\GeV$),
and Higgs ($125.25\GeV$) bosons alongside the 14~pseudoscalar mesons
adds 476 new triples (times sign patterns and both $Q(m)$ and
$Q(1/m)$---some~7800 trials in total).  The dual ratio
$Q(1/m) \equiv \sum 1/m_k \bigl/ \bigl(\sum 1/\sqrt{m_k}\bigr)^2$
is of interest because the Seiberg seesaw maps $m_j \to 1/m_j$.

\emph{Results:}
The best hit involving a gauge or Higgs boson is
$Q(1/m_{\pi^0},\, 1/m_{\eta_c},\, 1/m_W) = 0.6661$ ($0.08\%$
from~$2/3$, all-positive signs), but exact $Q = 2/3$ requires
$m_W = 82.6\GeV$ ($164\sigma$ from PDG)---excluded.
Among $Q(m)$ triples, $(\pi^0,\, B_c,\, H)$ reaches $0.15\%$
and $(D^\pm,\, D^0,\, H)$ reaches $0.21\%$, both all-positive;
but again exact $Q = 2/3$ requires $m_H \approx 126.5$--$126.9\GeV$
($7$--$10\sigma$ from PDG).  With $\sim 7800$~trials, roughly six
hits at the $0.1\%$ level are expected by chance.
No significant Koide relation exists among the electroweak bosons
or between bosons and mesons beyond what the look-elsewhere
correction explains.

%======================================================================
\section{The heavy triple: $(c, b, t)$}
%======================================================================

\subsection{A tuple without mesons}

The triple $(c, b, t)$ with all positive signs gives $Q = 0.6695$,
deviating by $0.42\%$. Its Koide scale $v_0^2 \approx 29.6\GeV$
is far above the QCD scale.

This triple has \textbf{no meson counterpart}: the top quark decays
weakly ($t \to Wb$, lifetime $\sim 5 \times 10^{-25}$~s) before the
strong interaction can bind it ($\tau_{\text{QCD}} \sim 1/\Lambda_{\text{QCD}}
\sim 10^{-24}$~s). No $T$-meson ($t\bar{q}$) exists. The meson matrix
$M^i{}_j$ extends only to $b$-flavored states; the top sector lives
purely at the quark level.

\subsection{The $(c,b,t)$ triple has no seed}

Unlike the $(0,s,c)$ and $(0,\pi,D_s)$ seeds, there is no obvious
$(0, x, y)$ seed for the $(c,b,t)$ triple. No pair of quarks in this
sector satisfies $\sqrt{m_{\text{light}}}/\sqrt{m_{\text{heavy}}}
\approx 2 - \sqrt{3}$: one has $\sqrt{m_c}/\sqrt{m_t} = 0.086$
(too small) and $\sqrt{m_c}/\sqrt{m_b} = 0.551$ (too large).

This means $(c,b,t)$ is \emph{not} a ``bloomed seed'' in the same
sense as $(-s,c,b)$. It may represent a fundamentally different
structure---a triple that was never seeded by a massless member but
instead emerged fully formed at a higher breaking scale, perhaps
connected to electroweak symmetry breaking rather than to chiral
symmetry breaking.

\subsection{Interpretations}

The $(c,b,t)$ triple ($Q = 0.6695$, $v_0^2 = 29.6\GeV$) has
neither a seed nor a meson shadow. It operates at a scale far
above $\Lambda_{\text{QCD}}$, suggesting its origin is different:

\begin{enumerate}
\item It may be governed by the electroweak Higgs mechanism rather
than the chiral condensate. Since $m_i = y_i\,v_{\text{EW}}/\sqrt{2}$, the Koide condition
$Q = 2/3$ translates into a constraint on the Yukawa couplings:
$\sum y_i/(\sum \sqrt{y_i})^2 = 2/3$.
With $y_t \approx 1$ fixed by the top mass, this determines
$y_b/y_t$ and $y_c/y_t$ through the Koide parametrization.

\item It may represent a Volkov--Akulov realization at the
electroweak scale: if there is a second SUSY breaking at
$f' \sim v_{\text{EW}}$, the $(c,b,t)$ triple would be its
analog of the $(\pi, \mu)$ system.

\item Or it may be a purely kinematic accident---the quark masses
happen to satisfy the Koide condition without a dynamical reason.
The 0.42\% precision makes this unlikely but not impossible.
\end{enumerate}

\subsection{CKM mixing and the Koide triples}

A striking feature of the quark Koide triples is that they
\emph{obligatorily mix up-type and down-type quarks}: $(-s,c,b)$
contains one up-type ($c$) and two down-type ($s,b$), while
$(c,b,t)$ contains two up-type ($c,t$) and one down-type ($b$). An
exhaustive scan of all $\binom{6}{3} = 20$ quark triples with all
sign choices confirms that \emph{every} triple close to $Q = 2/3$ is
mixed---neither the pure down-type $(d,s,b)$ nor the pure up-type
$(u,c,t)$ comes close ($Q = 0.73$ and $Q = 0.85$ respectively).
This mixing of mass eigenvalues from different Yukawa sectors is
puzzling in the Standard Model, where up-type and down-type masses are
eigenvalues of independent matrices $Y_u$ and $Y_d$. It is natural,
however, in an $SU(5)$ flavor-symmetric framework where all five
quarks are treated democratically.

The Koide constraints on the two overlapping triples, together with
the seed condition on $(0,s,c)$, determine $(s,c,b)$ from a single
input ($m_t$). The unconstrained masses are $m_u$ and $m_d$---the
first generation. The Gatto--Sartori--Tonin relation~\cite{GST1968}
$|V_{us}| \approx \sqrt{m_d/m_s}$ connects precisely these free
parameters to the Cabibbo angle ($|V_{us}| = 0.224$ vs.\ PDG $0.2243$),
suggesting that the CKM matrix
encodes the residual freedom left undetermined by the Koide
structure.

%======================================================================
\section{Mass relations and iterative chains}
\label{sec:chains}
%======================================================================

The preceding sections established the Koide energy-balance condition
$Q = 2/3$ and identified the triples that satisfy it. In this section
we quantify how special the lepton triple really is, introduce the
angular parametrization that reveals a hidden rational structure in
the Koide phase, and trace an iterative chain that connects the heavy
quarks to the strange sector. We report both successes and failures:
the scaling rule that generates the quark triple from the lepton
triple works once and then stops, and the iterative chain stalls at
a mass too light for any known fermion. These negative results
constrain the scope of the pattern.

\subsection{The lepton Koide formula: how special?}
Using PDG~2024 values $m_e = 0.51100\MeV$, $m_\mu = 105.658\MeV$,
$m_\tau = 1776.86 \pm 0.12\MeV$:
\begin{equation}
Q(e,\mu,\tau) = \frac{(\sqrt{m_e} + \sqrt{m_\mu} + \sqrt{m_\tau})^2}
{m_e + m_\mu + m_\tau} = 1.500014 \pm 0.000015\,,
\end{equation}
where $Q = (\sum\sqrt{m})^2 / \sum m$ is the reciprocal of the
convention in \S2: the Koide condition is
$Q = 3/2$ here, equivalently $\sum m / (\sum\sqrt{m})^2 = 2/3$
there. We use the $3/2$ convention in this section to match
the most common form in the Koide literature.
The deviation from $3/2$ is $|Q - 3/2| = 1.38 \times 10^{-5}$,
corresponding to $0.91\sigma$ given the uncertainty dominated by
$\delta m_\tau$.

\paragraph{Look-elsewhere effect.}
The Standard Model contains 9~charged fermions (3~leptons, 6~quarks).
An exhaustive scan of all $\binom{9}{3} = 84$ mass triplets yields
the following ranking.

\begin{center}
\begin{tabular}{clc}
\toprule
Rank & Triplet & $|Q - 3/2|$ \\
\midrule
1 & $(e, \mu, \tau)$ & $1.38 \times 10^{-5}$ \\
2 & $(c, b, t)$ & $6.2 \times 10^{-3}$ \\
\bottomrule
\end{tabular}
\end{center}

The gap between rank~1 and rank~2 is a factor of~450.

To assess the statistical weight of the lepton result, we performed a
Monte Carlo study: $10{,}000$ random mass sets of 9~fermions drawn
log-uniform on $[0.1, 200{,}000]\MeV$ were scanned for the best
triplet. The fraction with $|Q - 3/2|_{\min} \leq 1.38 \times 10^{-5}$
gives a look-elsewhere-corrected significance of $\mathbf{2.8\sigma}$.
This is suggestive but not conclusive: the lepton Koide formula is
unlikely to be a pure accident, but the evidence does not reach the
discovery threshold.

\paragraph{Extended scan with mesons.}
Enlarging the pool to 30~particles (9~fundamental fermions, $W$, $Z$,
$H$, $\fpi$, $v_{\text{EW}}$, and 16~pseudoscalar and vector mesons)
with all $2^3$~sign combinations per triple yields $\binom{30}{3} = 4060$
unique triples.  The lepton triple remains rank~1
($|Q-3/2| = 1.4\times 10^{-5}$, $0.001\%$), but $(c,b,t)$ drops to
rank~113 ($0.42\%$) and $(-s,c,b)$ to rank~324 ($1.24\%$): 112~meson-containing
triples are closer to $Q = 3/2$ than $(c,b,t)$.  Twenty-one triples
fall within $0.1\%$ of $3/2$, versus $\sim 12$~expected from a uniform
distribution---a ratio of $1.7$, not statistically remarkable.
\emph{Only $(e,\mu,\tau)$ involves three independent fundamental
parameters; the significance of quark Koide triples is diluted by the
composite-meson pool.}

\subsection{The Koide angle parametrization}
Any triple satisfying $Q = 3/2$ exactly can be written~\cite{Koide1983}
\begin{equation}
\sqrt{m_k} = \sqrt{M_0}\left(1 + \sqrt{2}\cos\!\left(
\frac{2\pi k}{3} + \delta_0\right)\right), \qquad k = 0,1,2\,,
\end{equation}
where $M_0 > 0$ is the scale parameter and $\delta_0 \in [0, 2\pi)$
is the Koide phase. This is not an approximation: any three
non-negative numbers whose $Q$-ratio equals $3/2$ are exactly
parametrized by some $(M_0, \delta_0)$.

Fitting to the charged lepton masses:
\begin{equation}
M_0 = 313.84\MeV\,, \qquad \delta_0 = 2.3166\;\text{rad}\,.
\end{equation}
The scale $M_0 \approx 314\MeV$ is the Koide scale $v_0^2$ of
\S2, now identified as the mass parameter of the
angular decomposition.

\paragraph{The phase is close to $2/9$.}
Unlike $Q = 2/3$, which is the algebraic consequence of the
energy-balance condition (\S2), the phase
$\delta_0$ is a free parameter that the energy balance does
not fix. Its value is an empirical datum, not a derived quantity.
The parametrization has a $2\pi/3$ periodicity
in $\delta_0$ (permuting the three masses). The physically
meaningful quantity is the reduced phase
\begin{equation}
\delta_0 \bmod \frac{2\pi}{3} = 0.22222\;\text{rad}\,.
\end{equation}
This is strikingly close to the fraction $2/9 = 0.22222\ldots$:
\begin{equation}
\delta_0 \bmod \frac{2\pi}{3} - \frac{2}{9} = +7.4 \times 10^{-6}\,,
\end{equation}
a residual of $33\,\text{ppm}$ relative to $2/9$
(using PDG 2024 central masses; within the $m_\tau$ uncertainty
of $0.12\MeV$, the residual vanishes at $m_\tau \approx 1{,}776.97\MeV$,
i.e.\ $0.9\sigma$ above central).

\paragraph{Statistical significance of the rational approximation.}
There are 266~reduced fractions $p/q$ with $q \leq 20$ and
$0 < p/q < 2\pi/3$. The probability that a random phase (uniform
on $[0, 2\pi/3)$) falls within $7.4 \times 10^{-6}$ of \emph{any}
of these 266 fractions gives a look-elsewhere-corrected significance
of $\mathbf{3.1\sigma}$. Evaluated specifically against the fraction
$2/9$ alone (i.e., without look-elsewhere correction), the
significance is $\mathbf{4.5\sigma}$.
The deviation from~$2/9$ is $0.9\sigma$ in the
$m_\tau$~uncertainty, consistent with an exact match.

\subsection{The scaling rule: $M_1 = 3M_0$, $\delta_1 = 3\delta_0$}

A natural question is whether other Koide triples can be obtained
from the lepton parameters by a simple scaling. The ansatz
\begin{equation}
M_1 = 3\,M_0 = 941.5\MeV\,, \qquad
\delta_1 = 3\,\delta_0 = 6.9497\;\text{rad}\,,
\end{equation}
inserted into the Koide parametrization, predicts a quark triple.
The following table compares predictions to PDG values.

\begin{center}
\begin{tabular}{lccc}
\toprule
Quark & Predicted (MeV) & PDG (MeV) & Deviation \\
\midrule
$s$ & 92.3 & 93.4 & $1.2\%$ \\
$c$ & 1360 & 1270 & $7.1\%$ \\
$b$ & 4197 & 4180 & $0.4\%$ \\
\bottomrule
\end{tabular}
\end{center}

The agreement for $s$ and $b$ is at the percent level; charm shows a
larger ($7.1\%$) discrepancy. For the physical masses,
$M_0(-s,c,b)/M_0(e,\mu,\tau) = 2.91$, not exactly $3$---the
$\times 3$ rule is approximate. By contrast, the vacuum-charge
doubling (\S4) predicts $m_b$ to $0.1\sigma$.
The predicted triple satisfies the
Koide condition with a negative square root for the strange quark:
\begin{equation}
Q(-\sqrt{m_s},\;\sqrt{m_c},\;\sqrt{m_b})
= \frac{(-\sqrt{m_s} + \sqrt{m_c} + \sqrt{m_b})^2}{m_s + m_c + m_b}
= \frac{3}{2}\,,
\end{equation}
holding to machine precision (residual $2.2 \times 10^{-16}$), since
the masses are generated from the parametrization which
enforces $Q = 3/2$ identically.

\paragraph{Renormalization caveat.}
The PDG quark masses used in the comparison are quoted at different
scales: $m_s(2\GeV)$, $m_c(m_c)$, $m_b(m_b)$, and $m_t$ as the pole
mass. The predicted masses emerge at a single internal scale set by
$M_1$. Running all masses to a common $\overline{\text{MS}}$ scale
would change the individual deviations but
not the structural point: the $\times 3$ scaling produces a
recognizable quark triple. QCD mass \emph{ratios} are RG invariants
at leading order, so the Koide condition itself is scheme-independent,
but the comparison of absolute masses to PDG values carries an
$O(\alpha_s)$ ambiguity.

\subsection{The second scaling fails}
Iterating the scaling rule:
\begin{equation}
M_2 = 9\,M_0 = 2825\MeV\,, \qquad
\delta_2 = 9\,\delta_0 = 20.849\;\text{rad}\,,
\end{equation}
gives predicted masses $(478,\; 92.3,\; 16{,}377)\MeV$. Only the
middle value echoes $m_s$. The lightest ($478\MeV$) does not match
any known quark mass, and the heaviest ($16.4\GeV$) falls between
$m_b$ and $m_t$ with no candidate. No SM triplet accommodates these
values.

\textbf{The scaling stops at one step.} This is a negative result
that constrains any attempt to build a recursive tower of Koide
triples: the multiplicative rule $M \to 3M$, $\delta \to 3\delta$
is not a general principle but a one-off relation between the lepton
and quark sectors. Whether this single success is accidental or
reflects a deeper connection (perhaps related to the three colors
of QCD, or to the three generations) remains open.

\subsection{Iterative chain: descending from $t$ and $b$}

An alternative approach to connecting the quark masses is purely
algebraic: given two known masses, solve the Koide equation
$Q = 3/2$ for the third. Starting from the top and bottom quarks:

\paragraph{Step~1: $Q(m_b, m_?, m_t) = 3/2$.}
The Koide equation with $(m_b, m_t)$ known is quadratic in
$\sqrt{m_?}$, yielding four branches (two signs for $\sqrt{m_?}$,
two roots of each quadratic). The solutions are:
\begin{equation}
m_? = \{1357,\; 60{,}280,\; 1.34 \times 10^6,\;
3.55 \times 10^6\}\MeV\,.
\end{equation}
Only the first branch gives a physically viable mass:
$m_? = 1357\MeV$, which we identify with the charm quark
(PDG: $m_c(m_c) = 1270\MeV$, deviation $6.8\%$).

\paragraph{Step~2: $Q(m_c, m_?, m_b) = 3/2$.}
With $m_c = 1357\MeV$ (the chain value, not the PDG value) and
$m_b = 4180\MeV$:
\begin{equation}
m_? = 92.17\MeV\,,
\end{equation}
which we identify with the strange quark. The negative square root
emerges naturally from the algebra---the solution selects the branch
with $\sqrt{m_?} < 0$ in the signed-charge framework, consistent
with the $(-s, c, b)$ structure of \S3.

\paragraph{Step~3: the chain stalls.}
Continuing from $(m_s, m_c)$ with $m_s = 92.17\MeV$:
\begin{equation}
m_? = 0.035\MeV\,.
\end{equation}
This is far too light for either $u$ ($2.16\MeV$) or $d$
($4.67\MeV$). The chain does not reach the first-generation quarks;
it stalls at a mass two orders of magnitude below $m_u$.

\paragraph{Step~4: cyclic echo.}
The stalled mass provides an algebraic consistency check.
Taking $m_{\text{stall}} = 0.035\MeV$ and $m_s = 92.17\MeV$, and
solving $Q(m_{\text{stall}}, m_?, m_s) = 3/2$:
\begin{equation}
m_? = 1357\MeV \quad \text{(to machine epsilon)}\,.
\end{equation}
The chain returns to its starting point: $t \to c \to s \to
\text{stall} \to c \to \cdots$. This is not a numerical coincidence
but an algebraic necessity: the Koide condition, combined with the
specific masses of $b$ and $t$, defines a closed orbit in
$(m_1, m_2, m_3)$ space. The cycle is:
\begin{multline*}
\cdots \;\to\; (b,t) \xrightarrow{Q=3/2} c
\;\to\; (c,b) \xrightarrow{Q=3/2} s \\
\;\to\; (s,c) \xrightarrow{Q=3/2} \text{stall}
\;\to\; (\text{stall}, s) \xrightarrow{Q=3/2} c
\;\to\; \cdots
\end{multline*}

The chain captures $c$, $b$, $s$ but misses $u$, $d$, and the
leptons entirely. The first generation is outside the algebraic
orbit seeded by the third generation.

\subsection{Cross-checks and master comparison}
Two independent routes lead to the quark masses: the
lepton-seeded scaling rule and the
top-bottom iterative chain.
The following table compares all predictions.

\begin{center}
\begin{tabular}{lcccc}
\toprule
Quark & Chain A & Chain B & PDG & Status \\
 & (from $t,b$) & (from leptons) & & \\
\midrule
$s$ & 92.2 & 92.3 & 93.4 & $1\text{--}1.2\%$ deviation \\
$c$ & 1357 & 1360 & 1270 & $7\text{--}7.1\%$ deviation \\
$b$ & (input) & 4197 & 4180 & $0.4\%$ deviation \\
$t$ & (input) & --- & 172\,760 & --- \\
stall & 0.035 & --- & --- & no SM match \\
$M_2$ set & --- & $(478,\,92,\,16377)$ & --- & no SM match \\
\bottomrule
\end{tabular}
\end{center}

The convergence of the two chains on $m_s \approx 92\MeV$ and
$m_c \approx 1357\text{--}1360\MeV$ is the central positive result.
Both routes, starting from completely different inputs (lepton masses
vs.\ heavy quark masses), produce the same strange and charm masses
to within $0.2\%$ of each other. Against PDG values, the strange
quark prediction is excellent ($\sim 1\%$), and the bottom quark
prediction is even better ($0.4\%$); the charm prediction shows a
persistent $\sim 7\%$ overshoot.

The negative results are equally informative:
\begin{itemize}
\item The $\times 3$ scaling does not iterate: the $\times 9$ step
produces no viable triple.
\item The iterative chain does not reach the first generation: it
stalls at $0.035\MeV$, two orders of magnitude below $m_u$.
\item The chain is algebraically cyclic: it loops through
$(c, s, \text{stall})$ indefinitely, never escaping to new masses.
\end{itemize}
These failures delineate the boundary of the Koide pattern: it
connects the second and third quark generations (and, via the
scaling rule, the leptons), but does not extend to the first
generation or beyond the Standard Model spectrum. Whatever
dynamics underlies the energy-balance condition
$\langle z_k^2\rangle = z_0^2$ operates within a limited domain.

\paragraph{The isospin seed.}
The chain above fails to reach the first generation through
individual quark masses, but the isospin combination
$m_u + m_d = 6.83\MeV$ provides the missing link.
The seed Koide ratio
\begin{equation}
Q(0,\,\sqrt{m_u + m_d},\,\sqrt{m_s}) = 0.6649
\end{equation}
deviates from $2/3$ by only $0.27\%$---comparable to
the $(0,s,c)$ seed at $0.35\%$.
Since $Q(0,a,b) = 2/3$ requires $b/a = (2+\sqrt{3})^2$,
the exact seed predicts $m_s/(m_u{+}m_d) = (2{+}\sqrt{3})^2 = 13.93$,
giving
\begin{equation}
m_s = (2+\sqrt{3})^2\,(m_u + m_d) = 95.1\MeV\,,
\end{equation}
a $0.36\sigma$ prediction from the PDG value $93.4 \pm 4.8\MeV$.
The lattice average $m_s/m_{ud} = 27.23 \pm 0.10$
(FLAG 2021, $N_f = 2{+}1{+}1$)
measures this ratio far more precisely than individual
quark masses; since $m_{ud} \equiv (m_u{+}m_d)/2$,
the prediction $2(2{+}\sqrt{3})^2 = 27.86$ exceeds the lattice
value by $6.3\sigma$.
This deviation---$2.3\%$ in the ratio, $0.27\%$ in $Q$---is
the bloom: the physical spectrum is shifted from the exact seed
$Q = 2/3$, just as $(-s,c,b)$ deviates by $1.2\%$ and $(c,b,t)$
by $0.4\%$ in~$Q$.
Combined with the $m_s$ chain
($m_s \to m_c \to m_b \to m_t$), the entire heavy quark
spectrum descends from the single isospin parameter
$m_u + m_d$---the same combination that enters the
Gell-Mann--Oakes--Renner relation
$m_\pi^2 = (m_u + m_d)\,|\qqbar|/\fpi^2$.
Since $m_s/(m_u{+}m_d)$ is an RG invariant (all three masses
share the same anomalous dimension), the isospin seed is
scheme-independent.

The isospin seed does not contradict the chain's failure to
reach the first generation: it reaches the \emph{sum}
$m_u + m_d$ but not $m_u$ and $m_d$ individually.
The individual splitting $m_d - m_u$ requires separate input
(connected to the Cabibbo angle via the GST relation).

\subsection{The dual Koide}

At the Seiberg vacuum with $B = \bar B = 0$, the diagonal meson VEVs
scale as $\langle M^j{}_j\rangle \propto 1/m_j$: heavy quarks produce
light mesons and vice versa. This seesaw prompts a dual question: does
the inverse-mass spectrum $\{1/m_i\}$ satisfy a Koide condition?
Evaluating $Q = \sum(1/m_i)/(\sum\sqrt{1/m_i})^2$ for the pure
down-type quarks $(d,s,b)$:
\begin{equation}
  Q(1/m_d,\, 1/m_s,\, 1/m_b) = 0.6652\,,
\end{equation}
a 0.22\% deviation from $2/3$---better than any direct quark Koide
triple. No other quark triple comes close under
the Seiberg seesaw (the next best, $(d,s,c)$, deviates by 4.2\%).
This ``dual Koide'' on $(d,s,b)$ (first reported here) connects
the UV mass hierarchy to the IR meson VEVs: the Seiberg constraint
maps the down-type quark masses into a near-Koide spectrum of
meson condensates. However, the
result is sensitive to $m_d$: within $\pm 1\sigma$ ($m_d = 4.67
\pm 0.48\MeV$), $Q$ ranges from $0.655$ to $0.676$, and the deviation
from $2/3$ varies from $0.2\%$ to $1.7\%$. The dual Koide is an
observation at the edge of the current $m_d$ uncertainty.
Setting $Q(1/m_d,1/m_s,1/m_b) = 2/3$ exactly predicts
$m_b = 4562\MeV$ ($9\%$ from PDG, $12.7\sigma$)---much worse
than $v_0$-doubling ($m_b = 4177\MeV$, $0.07\%$).  The $0.22\%$
closeness of $Q$ to $2/3$ does not translate to a precise mass
prediction because the Koide quadratic amplifies small $Q$
deviations into large mass deviations.

A simple algebraic identity connects the seed to the seesaw:
for any two masses, the Koide ratio
$Q(a,b) = (a+b)/(\sqrt{a}+\sqrt{b})^2$ satisfies
$Q(a,b) = Q(1/a,\,1/b)$ identically.
The seed $(0, m_s, m_c)$ sits at this self-dual point:
its Koide ratio is the same whether evaluated on quark masses or
on Seiberg-dual meson VEVs $\propto 1/m_i$.
A systematic scan over all $\binom{6}{3} = 20$ quark triplets from
$\{u,d,s,c,b,t\}$, allowing each entry to appear as $\sqrt{m_k}$
or $1/\sqrt{m_k}$ with independent signs (640 cases total), confirms
that $(1/m_d,1/m_s,1/m_b)$ achieves the globally best Koide
approximation among all quark triples with deviation $|Q-2/3| =
0.00144$, outperforming even $(c,b,t)$ at $0.00397$.  Ranks 2--5 in
the scan are occupied by mixed triples of the form $(\sqrt{m_s},
\sqrt{m_c}, \pm 1/\sqrt{m_X})$ with $X = b$ or $t$, all clustering
near $Q \approx 0.664$.  These mixed triples are moreover
dimensionally ill-defined: $\sqrt{m}$ carries units
$\MeV^{+1/2}$ while $1/\sqrt{m}$ carries $\MeV^{-1/2}$, so
their sum---and hence $Q$---is physically meaningless.
A direct unit-dependence test confirms this: $Q(\sqrt{m_s},\sqrt{m_c},1/\sqrt{m_b})$
equals $0.664$ when masses are expressed in MeV but $0.434$ in GeV,
while $Q(1/m_d,1/m_s,1/m_b)$ and $Q(m_c,m_b,m_t)$ are
unit-independent, as required of any physical ratio.
Equivalently, any mixed triple can be rendered dimensionally consistent
by interpreting the inverse entry as a seesawed mass
$m_3 = M^2/m_c$; setting $Q = 2/3$ then predicts $m_3$,
and the scale $M^2 = m_\mathrm{exp}\,m_\mathrm{pred}$ is merely
a derived label with no independent content.
Setting aside the mixed cases, the remaining near-Koide hits are not independent
coincidences either.  For any two masses $A < B$, the degenerate ratio
$Q(A,B,0) \equiv (A+B)/(\sqrt{A}+\sqrt{B})^2$ equals $2/3$ exactly
when $B/A = (2+\sqrt{3})^2 = 7+4\sqrt{3} \approx 13.93$.
The actual ratio $m_c/m_s = 1270/93.4 \approx 13.59$ is only $2.4\%$
below this threshold, giving $Q(m_s,m_c,0) = 0.664$.
Since $1/\sqrt{m_b}$ and $1/\sqrt{m_t}$ are negligible compared
to $\sqrt{m_s}+\sqrt{m_c} \approx 45\,\MeV^{1/2}$, appending any
heavy-quark inverse as a third entry merely perturbs $Q(m_s,m_c,0)$
by order $10^{-4}$.
Ranks 2--5 therefore all reduce to the near-seed position of the
$(s,c)$ pair and carry no additional Koide information.  The three
genuinely independent near-Koide quark triples are
$(-s,c,b)$, $(c,b,t)$, and $(1/d,1/s,1/b)$.

The structure becomes transparent when the six quarks are
grouped into three doublets
$(t,s)\times(u,b)\times(c,d)$,
each Koide triple selecting one element from each pair.
Including limits $m_u\!\to\! 0$ and $m_t\!\to\!\infty$,
there are $3\times 3\times 2 = 18$ base mass triples spanning
$2^3 = 8$ families.  Only four of the eight produce
$|Q - 2/3| < 1.5\%$:
%
\begin{center}
\small
\begin{tabular}{lllr}
\textbf{Family} & \textbf{Best triple} & \textbf{Type} & $|Q-\tfrac{2}{3}|$ \\
\hline
SBD & $(1/s,\,1/b,\,1/d)$     & all-inverse  & $0.22\%$ \\
SUC & $(s,\,0,\,c)$            & seed limit   & $0.35\%$ \\
TBC & $(t,\,b,\,c)$            & all-direct   & $0.60\%$ \\
SBC & $(-s,\,b,\,c)$           & signed       & $1.25\%$ \\
\end{tabular}
\end{center}
%
The remaining four families---TUC, TUD, TBD, SUD---deviate
by $4.6\%$ or more, with TUC worst at $27\%$.
The SUC seed is the pair Koide $Q(m_s,m_c,0) = 0.664$ already
derived above.
The doublets $(t,s)$, $(u,b)$, $(c,d)$ pair each
up-type quark with a down-type from a \emph{different} generation,
related by a cyclic permutation---the same
structure that appears in the CKM off-diagonal entries.

\paragraph{Charge decomposition and the seesaw.}
The Koide condition $Q=2/3$ is equivalent to the energy-balance
condition $\langle z_k^2\rangle = z_0^2$ (see \S2),
where $z_0$ is the mean charge and $z_k = \zch_k - z_0$ are the
perturbations.
The dual Koide on $(1/m_d,1/m_s,1/m_b)$ admits the
\emph{same} decomposition in terms of dual charges
$\xi_k = 1/\sqrt{m_k}$: writing $\xi_k = \xi_0 + \xi'_k$,
one finds $\langle\xi_k'^{\,2}\rangle/\xi_0^2 = 0.996$
($0.4\%$ from unity), closer to exact energy balance than
the direct triple $(-s,c,b)$ at $1.025$ ($2.5\%$).

However, the seesaw map $\zch \to 1/\zch$ does \emph{not}
preserve the Koide parametrization
$\zch_k = z_0(1+\sqrt{2}\cos(\delta+2\pi k/3))$.
Substituting $1/\zch_k = \xi_0(1+\sqrt{2}\cos(\delta'+2\pi k/3))$
and requiring the product $\zch_k \cdot 1/\zch_k = 1$
leads to a quadratic in $\cos\psi_k$ that must hold for
three distinct cosine values---impossible unless
one $\zch_k = 0$ (the seed).
The seed self-duality $Q(0,a,b) = Q(0,1/a,1/b)$ is
thus the \emph{unique} fixed point of the seesaw:
the electric and magnetic Koide conditions coincide
only at the SUSY-preserving vacuum where one mass vanishes.

The three independent Koide triples therefore
decompose into two sectors:
\begin{itemize}
\item \textbf{Electric} (charges $\zch = \pm\sqrt{m}$):
$(-s,c,b)$ and $(c,b,t)$, constraining $\{s,c,b,t\}$.
\item \textbf{Magnetic} (charges $\xi = 1/\sqrt{m}$):
$(1/d,1/s,1/b)$, constraining $\{d,s,b\}$.
\end{itemize}
The dual charges $\xi_j = \sqrt{\langle M^j{}_j\rangle}/\Lambda$
are literally the square roots of the Seiberg meson VEVs:
the magnetic Koide is an energy balance on the meson
condensate, not on quark masses.
Together with the $v_0$-doubling
($\sqrt{m_b} = 3\sqrt{m_s} + \sqrt{m_c}$),
these four constraints leave only $m_u$, $m_d$, and
the overall scale $m_s$ as free parameters.
If the isospin seed is included ($m_s = (2+\sqrt{3})^2(m_u+m_d)$,
\S\ref{sec:chains}), $m_s$ is predicted and
the magnetic Koide then determines $m_d$ from $(m_s,m_b)$
at $m_d = 4.61\MeV$ ($-0.13\sigma$);
the quark sector reduces to a single free parameter $m_u$.

Among all $\binom{5}{3} = 10$ quark triples, $(d,s,b)$ is the
unique triple with inverse-mass Koide ratio within 1\% of $2/3$;
the next closest is $(d,s,c)$ at 4.2\% deviation---a 19:1 margin.

\subsection{The energy balance is not an RG attractor}

One might hope that the energy-balance condition
$\langle z_k^2\rangle = z_0^2$ is a \emph{dynamical attractor}:
that three masses evolving under SUSY-breaking dynamics flow toward
the Koide surface $Q = 2/3$. However, the one-loop soft-mass RG
equations in $\mathcal{N}=1$ SQCD with universal Casimirs and
negligible Yukawas give $dm_i/dt = \gamma(M_g^2 - m_i)$, which
drives all masses to a common value $m_i \to M_g^2$, yielding
$Q \to 1/3$---the \emph{opposite} extreme. The Koide preservation
condition $\sum_i F_i(3 - 2S/\sqrt{m_i}) = 0$ on the $Q = 2/3$
surface is algebraically incompatible with any uniform attractor
flow (by the Cauchy--Schwarz inequality, the required identity
$S\sum 1/\sqrt{m_i} = 9/2$ contradicts $S\sum 1/\sqrt{m_i} \geq 9$).

The energy-balance condition therefore does \emph{not} arise from
simple RG running. It must instead be a \emph{boundary condition}
on the superpotential---as in the O'Raifeartaigh realization
discussed in~[I], where $gv = \sqrt{3}\,m$ produces the exact
Koide seed. The condition is set at the scale of SUSY breaking and
then preserved by the QCD-invariant mass ratios, rather than being
an RG attractor.

%======================================================================
\section{The de Vries angle and the Casimir equation}
%======================================================================

\subsection{The Casimir eigenvalue equation}

Consider the eigenvalue equation
\begin{equation}
x^2 \;+\; C_2\,x \;-\; C_2 = 0\,,
\qquad C_2 = s(s+1)\,,
\end{equation}
where $C_2$ is the quadratic Casimir of $SU(2)$ at spin $s$ and
$x = M^2/m^2$ for a mass scale $m$ to be determined.  Equivalently,
$x^2/(1-x) = C_2$: the eigenvalue is a self-consistent solution
where the ``squared coupling'' $x^2$ equals $C_2$ times the
complement $(1-x)$.  The solutions are
\begin{equation}
x_{\pm}(s) = \frac{-C_2 \pm \sqrt{C_2^2 + 4C_2}}{2}\,.
\end{equation}

\paragraph{Structural constraints (codimension analysis).}
The most general two-parameter family of such equations is
\begin{equation}
x^2 + (f_0 + f_1 C_2)\,x + (g_0 + g_1 C_2) = 0\,,
\end{equation}
with four real coefficients $(f_0, f_1, g_0, g_1)$.  Three
structural conditions pin the polynomial to the equation above:
\begin{enumerate}
\item $f_0 = 0$: no spin-independent mass term;
\item $g_0 = 0$: a massless eigenstate exists at $s = 0$;
\item $g_1 = -f_1$: coefficient antisymmetry
  (the constant and linear terms in $C_2$ are related).
\end{enumerate}
Together with the normalization $f_1 = 1$, these exhaust the
freedom.  A systematic scan of 582 polynomial alternatives
(quadratics, cubics, quartics with Casimir-type functions and
rational coefficients) confirms uniqueness: no other
equation of this class reproduces $\sin^2\!\theta_W$ to
sub-percent accuracy at $(s_1, s_2) = (1/2, 1)$.
The ratio derived below is therefore a
\emph{consequence} of these constraints, not a free parameter.

\subsection{The electroweak ratio}
Define the ratio
\begin{equation}
R \;\equiv\; 1 - \frac{x_+(1/2)}{x_+(1)}\,.
\end{equation}
Evaluating at $s = 1/2$ ($C_2 = 3/4$)
and $s = 1$ ($C_2 = 2$):
\begin{align}
x_+(1/2) &= \tfrac{1}{8}\bigl(-3 + \sqrt{57}\bigr)
  \,,\quad\text{so } x_+(1/2)\,m^2 = M_W^2\,,
\\[4pt]
x_+(1) &= -1 + \sqrt{3}
  \,,\quad\text{so } x_+(1)\,m^2 = M_Z^2\,.
\end{align}
The identification with $W$ and $Z$ is motivated by their
$SU(2)$ quantum numbers ($s = 1/2$ maps to the doublet structure of $W$,
$s = 1$ to the triplet).  Then
\begin{equation}
\boxed{R = 1 - \frac{x_+(1/2)}{x_+(1)}
  = \frac{(\sqrt{19} - 3)(\sqrt{19} - \sqrt{3}\,)}{16}
  \;\approx\; 0.2231013\ldots}
\end{equation}
The factored form follows from the discriminants:
$57 = 3 \times 19$ at $s = 1/2$ and
$12 = 4 \times 3$ at $s = 1$.  Rationalizing the denominator
$\sqrt{3} - 1$ yields an expression in $\sqrt{19}$ and $\sqrt{3}$
alone, which coincides term-by-term with the form found
empirically by de~Vries.

This should be compared with the on-shell electroweak mixing
parameter:
\begin{equation}
\sin^2\!\theta_W^{\text{os}}
  \;\equiv\; 1 - \frac{M_W^2}{M_Z^2}
  = 0.22320 \pm 0.00026\,,
\end{equation}
using $M_Z = 91.1876 \pm 0.0021\GeV$ and
$M_W = 80.369 \pm 0.013\GeV$ (PDG~2024).

Since $R = 1 - x_+(1/2)/x_+(1)$ and
$\sin^2\!\theta_W = 1 - M_W^2/M_Z^2$, the two expressions are
\emph{algebraically identical} once one identifies
$x_+(s) = M^2(s)/m^2$.  The prediction $R \approx
\sin^2\!\theta_W$ and the $W$-mass prediction derived below are
therefore \textbf{one and the same observation}, not two
independent tests.

\paragraph{$W$-mass prediction.}
From $M_W = M_Z\sqrt{1 - R}$:
\begin{equation}
M_W^{\text{pred}}
  = 91.1876 \times \sqrt{1 - 0.22310}
  = 80.374 \pm 0.002\GeV\,,
\end{equation}
which lies $0.39\sigma$ from the PDG average
$M_W = 80.369 \pm 0.013\GeV$ and strongly disfavors the
anomalous CDF-II measurement $M_W^{\text{CDF}} = 80.4335 \pm
0.0094\GeV$ (6.2$\sigma$ tension with $R$).

\paragraph{Scale parameter.}
Setting $x_+(1)\,m^2 = M_Z^2$ gives
\begin{equation}
m = \frac{M_Z}{\sqrt{\sqrt{3} - 1}}
  = 106.577\GeV
  \approx m_\mu \times 10^3\,,
\end{equation}
a coincidence whose significance is unclear.

\subsection{Open issues for the Casimir equation}

\paragraph{Renormalization scheme dependence.}
The comparison uses the on-shell
definition $\sin^2\!\theta_W = 1 - M_W^2/M_Z^2 = 0.22320$,
\textbf{not} the $\overline{\text{MS}}$ value
$\sin^2\!\theta_W(\overline{\text{MS}}) = 0.23122 \pm 0.00004$.
The two differ by $\Delta \approx 0.008$, far larger than
experimental errors.  The algebraic argument
contains no dynamical mechanism to
prefer one scheme over the other.  Until such a mechanism is
identified, the on-shell agreement should be regarded as
suggestive but scheme-dependent.

\paragraph{The fourth eigenvalue.}
At $s = 1/2$, the second root gives
\begin{equation}
|M(1/2, -)|
  = m\,\sqrt{|x_-(1/2)|}
  = 122.4\GeV\,,
\end{equation}
which is 2.3\% below $M_H = 125.25 \pm 0.17\GeV$.  However,
this eigenstate carries spin $1/2$, while the Higgs boson is
spin $0$.  The near-degeneracy is noted as a numerical curiosity;
we do not claim a connection.

\paragraph{Origin.}
This observation was first made by H.~de~Vries in a Physics
Forums discussion (November 2004) and subsequently studied in~\cite{Rivero2005b}.

\subsection{The Gatto--Sartori--Tonin relation}

A further mass relation connected to the Koide web is the GST
relation~\cite{GST1968}:
\begin{equation}
\sin\theta_C \;=\; \sqrt{\frac{m_d}{m_s}} \;=\; 0.2236\,,
\end{equation}
to be compared with the PDG value $|V_{us}| = 0.2243 \pm 0.0005$.
The deviation is $0.3\%$ ($1.4\sigma$).

This is not a Koide triple, but it connects to the Koide structure
through the free parameters: after the Koide seed fixes $m_c/m_s$,
the $v_0$-doubling fixes $m_b$, and the $(c,b,t)$ triple
overdetermines $m_t$, the remaining free masses are $m_u$ and $m_d$.
The GST relation connects $m_d/m_s$ to the Cabibbo angle, so the
CKM matrix encodes exactly the freedom left over after the Koide
conditions are imposed.

%======================================================================
\section{Statistical assessment}
%======================================================================

We now compile the statistical significance of each observation,
distinguishing between individually significant relations and those
that derive their weight from the structural context.

\subsection{Summary of individual significances}

\begin{center}
\small
\begin{tabular}{llccc}
\toprule
Observation & Precision & Raw $\sigma$ & LEE trials & LEE $\sigma$ \\
\midrule
$Q(e,\mu,\tau) = 2/3$ & $0.001\%$ & $0.91\sigma$ (in $m_\tau$) & 84 triples & $2.8\sigma$ \\
$\delta_0 \bmod 2\pi/3 = 2/9$ & 33 ppm & $0.9\sigma$ (in $m_\tau$) & 266 fractions & $3.1\sigma$ \\
$Q(-\pi,D_s,B) = 2/3$ & $0.10\%$ & --- & 660 combinations & $1.1\sigma$ \\
$Q(c,b,t) = 2/3$ & $0.42\%$ & --- & 84 triples & $<1\sigma$ \\
$Q(-s,c,b) = 2/3$ & $1.24\%$ & --- & 84 triples & $<1\sigma$ \\
$v_0$-doubling: $m_b = 4177$ & $0.07\%$ & $0.1\sigma$ & --- & not significant alone \\
Dual Koide: $Q(1/m_d,\ldots)$ & $0.22\%$ & --- & 10 triples & not significant alone \\
$R = \sin^2\theta_W$ & $0.04\%$ & $0.13\sigma$ (in $M_W$) & unique & --- \\
GST: $\sin\theta_C = \sqrt{m_d/m_s}$ & $0.3\%$ & $1.4\sigma$ & --- & --- \\
\bottomrule
\end{tabular}
\end{center}

\subsection{Which observations are individually significant?}

Two observations pass the $3\sigma$ threshold after look-elsewhere
correction:

\begin{enumerate}
\item \textbf{The lepton Koide formula} ($2.8\sigma$). Among 84 SM
fermion triples, the lepton triple is the best by a factor of 450
in $|Q - 2/3|$. A Monte Carlo with $10^4$ random mass spectra
confirms the $0.24\%$ false-positive rate.

\item \textbf{The phase relation $\delta_0 \bmod 2\pi/3 = 2/9$}
($3.1\sigma$ with LEE, $4.5\sigma$ specific). Among 266 reduced
fractions with denominator $\leq 20$, the phase falls within
$7.4 \times 10^{-6}$ of $2/9$.
\end{enumerate}

The de Vries angle ($R = 0.22310$ vs.\ $\sin^2\theta_W = 0.22320$)
gives $0.13\sigma$ in $M_W$, but assigning a look-elsewhere
correction is difficult because the trial space of ``algebraic
equations involving Casimir operators'' is not well defined.

\subsection{Which observations are not individually significant?}

The remaining observations derive their interest from context:

\begin{itemize}
\item \textbf{The $v_0$-doubling} ($m_b = 4177\MeV$, $0.1\sigma$
from PDG). In isolation, the look-elsewhere probability is $24.9\%$:
there are many linear relations among square roots of masses, and
some will match by chance. The significance comes from its role
within the seed-to-bloom framework.

\item \textbf{The meson triple} ($1.1\sigma$ after LEE). Not
significant on its own, but the structural pattern (all best
matches involve $-\pi$ with heavy mesons) is noteworthy.

\item \textbf{The quark triples} $(-s,c,b)$ and $(c,b,t)$.
Neither reaches $1\sigma$ after LEE. Their significance is that
they chain together and share masses with the lepton-seeded
predictions.

\item \textbf{The dual Koide} ($Q = 0.665$, $0.22\%$). Sensitive
to $m_d$ uncertainty and not predictive for $m_b$. Its interest is
the connection to the Seiberg seesaw.
\end{itemize}

\subsection{The correlations argument}

The key observation is that these relations are not independent.
They share masses, share algebraic structure, and point in the
same direction:

\begin{enumerate}
\item The lepton Koide, the $\times 3$ scaling, and the iterative
chain all converge on $m_s \approx 92\MeV$ and $m_c \approx 1357\MeV$.

\item The seed condition, the $v_0$-doubling, and the signed triple
all involve the same three masses $(m_s, m_c, m_b)$.

\item Every quark triple near $Q = 2/3$ obligatorily mixes up-type
and down-type quarks.

\item The Koide phase $\delta_0 = 2/9$ (mod $2\pi/3$) and the
$\times 3$ scaling $\delta_1 = 3\delta_0$ use the same parameter.

\item The free parameters after Koide ($m_u$, $m_d$) are connected
to the Cabibbo angle via GST.
\end{enumerate}

No formal combination of $p$-values is attempted, because the
correlations invalidate any naive product of independent
probabilities. Instead, we note that the \emph{pattern} of
interlocking relations---rather than any single observation---constitutes
the evidence. The probability of the full pattern arising by
chance is much lower than the product of individual $p$-values
would suggest, precisely because the observations are
\emph{structurally} related rather than statistically independent.

%======================================================================
\section{The electron mass and symmetry restoration}
%======================================================================

\subsection{The electron mass and the lepton seed}

The charged lepton triple $(e, \mu, \tau)$ is the ``almost-seed'':
$\zch_e = 0.715\MeV^{1/2}$ is small but nonzero. The electron's
charge is near zero because of cancellation in $\zch_e = z_0 + z_e$:
\begin{equation}
v_0 = 17.716, \quad z_e = -17.001, \quad
\zch_e = 17.716 - 17.001 = 0.715\MeV^{1/2}\,.
\end{equation}
A 96\% cancellation. The mass $m_e = 0.715^2 = 0.511\MeV$ is the
square of something already small.

The near-vanishing of $m_e$ is a structural property of the Koide
embedding, not fine-tuning: the energy balance allows perturbations
as large as $z_0$. Every Koide triple has one member with
$|\zch|/v_0$ small:

\begin{center}
\begin{tabular}{lccl}
\toprule
Triple & Lightest & $|\zch|/v_0$ & Character \\
\midrule
$(e, \mu, \tau)$ & $e$ & 0.040 & Almost-seed \\
$(-\pi, D_s, B)$ & $\pi$ (neg.) & 0.337 & Full triple \\
$(c, b, t)$ & $c$ & 0.236 & Full triple \\
$(-s, c, b)$ & $s$ (neg.) & 0.320 & Full triple \\
\bottomrule
\end{tabular}
\end{center}

The electron is by far the most extreme case, consistent with the
lepton triple being nearest to its seed.

\subsection{Symmetry restoration}

The broken spectrum reassembles as symmetry-breaking effects are
removed:

\begin{enumerate}
\item \textbf{Turn off QED:} Electromagnetic splittings vanish
($\pi^\pm = \pi^0$, $K^\pm = K^0$). D-term effects disappear.
\item \textbf{Set $m_u = m_d$:} Exact isospin.
\item \textbf{Set $m_s = m_d$:} $SU(3)_V$ restored,
$\pi = K = \eta$. Koide splitting $v_1 \to 0$; lepton partners
degenerate.
\item \textbf{Send $m_q \to 0$:} Chiral limit, all meson charges
$\to 0$. If sBootstrap SUSY restores, all light lepton masses
vanish too. This is the seed configuration: $(0, 0, 0)$.
\end{enumerate}

The seed triples are the ``penultimate'' stage: one member has
reached zero, but the other two retain the $2 - \sqrt{3}$ ratio.
Full triples are the final broken state. The lepton triple,
trapped between seed and bloom, reveals the hierarchy:
chiral breaking $\to$ seed $\to$ bloom $\to$ full spectrum.

The direction of the flow also matters: approaching the chiral
limit, the pion charge rises from $\zch = -11.81$ toward zero
from below. It crosses zero exactly in the chiral limit. This
means the pion was ``born negative''---its charge sign is set by
the direction of chiral symmetry breaking, and the seed
$(0, \pi, D_s)$ is the configuration at the moment of crossing.

%======================================================================
\section{Conclusions}
%======================================================================

The principal results of this paper are:

\paragraph{1.\ The Koide formula is energy equipartition.}
$Q = 2/3$ is $\langle z_k^2\rangle = z_0^2$: perturbation energies
match the base scale. The signed-charge framework resolves the
electron mass puzzle (near-cancellation, not fine-tuning) and reveals
that mesons satisfy Koide with $(-\pi, D_s, B)$ to $0.1\%$.

\paragraph{2.\ Seed triples and the $v_0$-doubling.}
The quark seed $(0,s,c)$ and meson seed $(0,\pi,D_s)$ share the
universal ratio $2 - \sqrt{3}$ and operate at $v_0^2 \sim
230$--$350\MeV$, close to the lepton triple's $v_0^2 = 314\MeV
\approx \fpi^2/27$. The bloom from seed to full triple doubles the
vacuum charge, yielding $\sqrt{m_b} = 3\sqrt{m_s} + \sqrt{m_c}$
($m_b = 4177\MeV$, $0.1\sigma$ from PDG)---more precise than the
overlapping-triples prediction.

\paragraph{3.\ The lepton Koide and the phase relation are
individually significant.}
$Q(e,\mu,\tau) = 2/3$ to $0.001\%$ gives $2.8\sigma$ after look-elsewhere;
$\delta_0 \bmod 2\pi/3 = 2/9$ to 33~ppm gives $3.1\sigma$ with LEE
($4.5\sigma$ specific). These are the strongest individual results.

\paragraph{4.\ The quark Koide triples obligatorily mix up-type
and down-type quarks.}
An exhaustive scan confirms that \emph{every} triple near $Q = 2/3$
mixes charge sectors. The pure same-charge triples $(d,s,b)$ and
$(u,c,t)$ give $Q = 0.73$ and $Q = 0.85$, both far from $2/3$.
The free parameters after Koide are $m_u$ and $m_d$, connected to
the Cabibbo angle via the GST relation.

\paragraph{5.\ An algebraic electroweak ratio.}
The Casimir eigenvalue equation $x^2 + s(s+1)\,x - s(s+1) = 0$
predicts $R = 1 - x_+(1/2)/x_+(1) = 0.22310$, matching the on-shell
$\sin^2\theta_W = 0.22320$ to $0.4\sigma$ and giving
$M_W = 80.374\GeV$ ($0.39\sigma$ from PDG, $6.2\sigma$ from CDF-II).
The equation is pinned by three structural constraints.

\paragraph{6.\ Negative results constrain the pattern.}
The $\times 3$ scaling does not iterate ($\times 9$ fails);
the iterative chain stalls at $0.035\MeV$; the Koide condition
is not an RG attractor (one-loop flow gives $Q \to 1/3$);
the cyclic echo is algebraically exact.
These failures delineate the domain of the energy-balance condition.

\paragraph{7.\ The energy-balance condition is a boundary condition,
not numerology.}
The pattern of mass relations compiled here---lepton Koide,
quark chains, meson triple, $v_0$-doubling, phase relation, de Vries
angle, and their interlocking structure---is not explained by any
known Standard Model mechanism. Nor is it dismissed by look-elsewhere
corrections: the two individually significant results ($2.8\sigma$
and $3.1\sigma$) and the correlated web of relations argue for a
common underlying principle.

The companion paper~[I] identifies this principle: the O'Raifeartaigh
superpotential with $gv/m = \sqrt{3}$ produces the Koide seed
exactly, and the layered SUSY-breaking structure generates the full
spectrum. The present paper provides the empirical foundation; [I]
provides the dynamical explanation.

\subsection{Epistemological classification}

The following table classifies this paper's claims by their logical
status.

\begin{center}
\small
\begin{tabular}{lll}
\toprule
Claim & Status & Comment \\
\midrule
$Q = 2/3$ iff $\langle z_k^2\rangle = z_0^2$ & Derived & Algebra \\
Seed ratio $= 2-\sqrt{3}$ & Derived & From $Q = 2/3$ + one zero \\
Cyclic echo ($t \to c \to s \to \text{stall} \to c$) & Derived & Algebraic necessity \\
Seed self-dual: $Q(a,b) = Q(1/a,1/b)$ & Derived & Identity for two masses \\
\midrule
Lepton Koide $Q = 2/3$ to $0.001\%$ & Observed & $2.8\sigma$ after LEE \\
Meson triple $(-\pi,D_s,B)$ to $0.10\%$ & Observed & $1.1\sigma$ after LEE \\
$v_0$-doubling: $v_0^{(f)}/v_0^{(s)} = 2.0005$; $m_b$ to $0.07\%$ & Observed & Derived via bion K\"ahler in~[I] \\
$\delta_\ell \bmod 2\pi/3 = 2/9$ to 33~ppm & Observed & $3.1\sigma$ after LEE \\
Casimir $R = 0.22310$ vs $\sin^2\theta_W$ & Observed & Scheme-dependent \\
Dual Koide $Q(1/m_d,1/m_s,1/m_b) = 0.665$ & Observed & $0.2\%$ from $2/3$; not predictive \\
GST: $|V_{us}| = \sqrt{m_d/m_s} = 0.224$ & Observed & $0.3\%$ from PDG \\
$M_0 \approx \fpi^2/27$ & Observed & $0.04\%$; connects lepton scale to $\fpi$ \\
\bottomrule
\end{tabular}
\end{center}

\begin{thebibliography}{99}

\bibitem{Koide1983}
Y.~Koide, ``New viewpoint of quark and lepton mass hierarchy,''
Phys.\ Rev.\ D \textbf{28}, 252 (1983).

\bibitem{PDG2024}
Particle Data Group, R.L.~Workman et~al.,
Prog.\ Theor.\ Exp.\ Phys.\ \textbf{2022}, 083C01 (2022);
updated 2024.

\bibitem{Foot1994}
R.~Foot, arXiv:hep-ph/9402242.

\bibitem{Koide1990}
Y.~Koide, Mod.\ Phys.\ Lett.\ A \textbf{5}, 2319 (1990).

\bibitem{Ma2006}
E.~Ma, Phys.\ Rev.\ D \textbf{73}, 077301 (2006).

\bibitem{XZZ2008}
Z.-z.~Xing, H.~Zhang, and S.~Zhou,
Phys.\ Lett.\ B \textbf{666}, 462 (2008).

\bibitem{Esposito1995}
S.~Esposito and P.~Santorelli,
Mod.\ Phys.\ Lett.\ A \textbf{10}, 3077 (1995).

\bibitem{RodejohannZhang2011}
W.~Rodejohann and H.~Zhang,
Phys.\ Lett.\ B \textbf{698}, 152 (2011).

\bibitem{Rivero2012}
A.~Rivero, ``A new Koide tuple: strange, charm, bottom,''
arXiv:1111.7232 [hep-ph] (2011).

\bibitem{Rivero2005}
A.~Rivero, ``An interpretation of scalars in SO(32),''
Eur.\ Phys.\ J.\ C \textbf{84}, 1058 (2024),
arXiv:2407.05397 [hep-ph].

\bibitem{Rivero2007}
A.~Rivero, arXiv:0710.1526 [hep-ph] (2007).

\bibitem{Zenczykowski2012}
P.~Zenczykowski, Phys.\ Rev.\ D \textbf{86}, 117303 (2012).

\bibitem{Rivero2005b}
A.~Rivero, ``The de Vries formula for $\sin^2\theta_W$,''
arXiv:hep-ph/0510068 (2005).

\bibitem{GST1968}
R.~Gatto, G.~Sartori, and M.~Tonin,
``Weak self-masses, Cabibbo angle, and broken $SU(3)$,''
Phys.\ Lett.\ B \textbf{28}, 128 (1968).

\bibitem{Harari}
H.~Harari, H.~Haut, J.~Weyers,
``Quark masses and Cabibbo angles,''
Phys.\ Lett.\ B \textbf{78}, 459 (1978).

\bibitem{Seiberg1995}
N.~Seiberg, ``Electric--magnetic duality in supersymmetric
non-abelian gauge theories,'' Nucl.\ Phys.\ B \textbf{435}, 129
(1995).

\end{thebibliography}

\tableofcontents

\end{document}
